\section{讨论与局限}\label{sec:limits}

\subsection{结论边界声明}

为使本文主张的范围明确，我们在此集中声明四条边界。这四条不是致歉式的免责，而是对“本文能支持什么、不能支持什么”的精确刻画。

\begin{enumerate}[leftmargin=2.2em, itemsep=2pt]
  \item \textbf{不提出新模型，也不涉及训练。}本文的三层模型、失效分类、验证梯度与核查清单是对现有流程的\emph{结构化描述}与\emph{可测化重述}，不包含任何新的网络结构、训练目标或参数规模的实验。因此本文的贡献是概念与工具层面的，不能据此推断某个具体模型的性能。
  \item \textbf{四类失效模式不主张完备性。}F1--F4 的划分依据是可操作性与可修复性，不是穷尽枚举。存在不落入这四类的失效（例如涉及数值精度、模型对符号约定的地域性差异等），本文未覆盖。任何“未检出 F1--F4 即认为规约正确”的使用方式都超出了本文的支持范围。
  \item \textbf{不作大样本经验主张。}$V$-rate、$\mathrm{Fid}$、$\mathrm{CDR}$ 三个指标的定义与可计算性不依赖样本规模，但本文\emph{不提供}在特定数据集上的测量结果，也不主张某一层次或某一指标在总体上优于其他。第 \ref{sec:case} 节的两个例子是为说明机制而构造的教学实例，不构成经验证据。
  \item \textbf{不替代形式化与自动形式化路线。}本文的框架与 Lean/mathlib 等路线是互补关系：形式化路线提供最强的 L3，本文关注的是进入 L3 之前的规约质量与归因链条。我们既不主张清单能替代内核验证，也不主张内核验证不需要规约检查。
\end{enumerate}

\subsection{局限与未解决问题}

\paragraph{$\mathrm{Fid}$ 的判据依赖标准规约。}在缺少权威形式化版本的题目上，$\mathrm{Fid}$ 的“正确规约”本身有争议，此时该指标退化为对标注者判断的一致性度量，而非对客观正确性的度量。

\paragraph{清单的执行成本。}完整档十项的执行依赖人工判断，尤其第 6 项（反译可读性）需要读者具备把形式表示还原为自然语言的能力。清单目前未实现自动化；把它转化为可半自动执行的检查（例如由检查器辅助逐项提问）是后续工作。

\paragraph{验证层级的标注主观性。}$V$-rate 依赖对每题实际达到层级的判定，边界情形（例如“用另一次模型判断但这判断本身被结构化约束”）在 V1 与 V2 之间存在归属争议。可行的缓解是要求标注同时给出依据，但这不消除争议。

\paragraph{语义漂移的自动检测。}F4 的检测目前依赖清单第 7 项的人工核对。能否通过追踪“每一步所依赖假设的最小集”来机械地检测假设被静默加强，是一个具体且值得尝试的方向；本文未给出算法。

\subsection{与既有路线的分工}

把本文的位置放在既有工作图谱上：生成式推理路线\citep{wei2022chain,kojima2022large}解决“如何产出推理”，其可靠性评估主要停在 V0/V1；工具增强与神经符号路线\citep{gao2022pal,chen2022program,schick2023toolformer,zhou2023solving}把计算外包给可执行环境，从而在\emph{计算}环节获得机械保证，但计算之外的自然语言到符号的转写仍是人工或模型判断；形式化与自动形式化路线\citep{moura2021lean,mathlib2020lean,wu2022autoformalization}提供最强的 L3，其瓶颈恰在 L2；过程级验证与可信评测路线\citep{lightman2023let,uesato2022solving,ni2023evaluating,biderman2024lessons}分别从步骤打分与评测效度两侧逼近可靠性。本文的差异化位置在于：把 L2 显式化为\emph{带类型与假设的中间表示}，并给出一组不依赖大规模实验即可计算的指标，使“验证做到了哪一层”成为一个可标注、可比较、可追责的量。

\section{结论}\label{sec:conclusion}

本文围绕一个观察展开：AI 数学系统在答案正确率上的提升，与结论的可验证性之间存在结构性错位。我们把这一错位定位到语义规约层，并给出四件互相支撑的产物：一个三层可靠性模型（生成层 L1、语义规约层 L2、验证层 L3）与三条由此导出的可检验推论；一个含四类失效模式（歧义未消解、类型不完整、不可逆、语义漂移）的分类；一个按\emph{验证器可访问的信息}划分的验证梯度 V0--V3 与之配套的三个可测指标（$V$-rate、$\mathrm{Fid}$、$\mathrm{CDR}$）；以及一份十项核查清单，其每一项都指向某个具体失效模式或指标。

核心论断可以用一句话概括：\emph{在语义规约层引入的误差对验证层不可见，因此仅增强验证层无法把可靠性提升到 $1-p_{\text{err}}$ 以上}（推论 R1）。第 \ref{sec:case} 节的两个例子示范了这一论断的具体形态——结论正确而推理断裂、结论正确而适用条件丢失，两种情形都逃过答案级比对，都必须由结构化检查捕获。

后续工作有三个具体方向：其一，把十项清单中的可判定项（第 3--5、7、8 项）做成半自动检查器，并测量其对人工发现的增量覆盖；其二，针对 F4 给出基于“步骤依赖假设最小集”的机械检测算法；其三，在具备权威形式化版本的题目上测量 $\mathrm{Fid}$ 与 $\mathrm{CDR}$ 的可复现性，从而检验两个指标在跨标注者条件下的稳定性。

\section*{代码与数据可用性}
本文的引文核验流程以脚本形式给出：给定本稿的 \texttt{references.bib}，脚本逐条向 arXiv API 与 Crossref API 查询，比对标题相似度并记录实际命中的来源 URL；相似度低于阈值的条目被列为未核验而不进入最终文献表。核验记录（含每条文献的来源 URL、命中方式与相似度）与脚本一并置于仓库的 \texttt{pipeline/} 与 \texttt{lit/} 目录。本文正文中的两个案例为自造教学实例，不含外部数据集。
